\begin{tikzpicture}[x=0.046cm,y=3.0cm]
% ============ top panel: xi_eq(dV) at three temperatures — eq:xieq, w_j=n R T/F, n=1 ============
\draw[-{Latex[length=1.6mm]}] (-108,0) -- (130,0) node[below,font=\scriptsize] {$V-U_j^{\,d}$ [mV]};
\draw[-{Latex[length=1.6mm]}] (-108,0) -- (-108,1.18) node[above,font=\scriptsize] {$\xi_\eq$};
\foreach \xx in {-100,-50,50,100} {\draw (\xx,0.02) -- (\xx,-0.02) node[below,font=\scriptsize] {\xx};}
\draw (-104,1)--(-112,1) node[left,font=\scriptsize] {1};
\draw (-104,0.5)--(-112,0.5) node[left,font=\scriptsize] {0.5};
\draw[densely dotted,black!45] (0,0) -- (0,1.02);
\draw[densely dashed] plot[smooth] coordinates {(-100,0.013) (-90,0.0199) (-80,0.0304) (-70,0.046) (-60,0.0693) (-50,0.1029) (-40,0.1503) (-30,0.2143) (-20,0.2961) (-10,0.3934) (0,0.5) (10,0.6066) (20,0.7039) (30,0.7857) (40,0.8497) (50,0.8971) (60,0.9307) (70,0.954) (80,0.9696) (90,0.9801) (100,0.987)};
\draw[thick] plot[smooth] coordinates {(-100,0.02) (-90,0.0292) (-80,0.0425) (-70,0.0615) (-60,0.0881) (-50,0.1249) (-40,0.174) (-30,0.2372) (-20,0.3146) (-10,0.4039) (0,0.5) (10,0.5961) (20,0.6854) (30,0.7628) (40,0.826) (50,0.8751) (60,0.9119) (70,0.9385) (80,0.9575) (90,0.9708) (100,0.98)};
\draw[densely dotted,very thick] plot[smooth] coordinates {(-100,0.0283) (-90,0.0398) (-80,0.0557) (-70,0.0775) (-60,0.1069) (-50,0.1457) (-40,0.1954) (-30,0.257) (-20,0.3301) (-10,0.4125) (0,0.5) (10,0.5875) (20,0.6699) (30,0.743) (40,0.8046) (50,0.8543) (60,0.8931) (70,0.9225) (80,0.9443) (90,0.9602) (100,0.9717)};
\node[right,font=\tiny,align=left] at (55,0.30) {$268$ K ($w{=}23.1$ mV, dashed)\\$298$ K ($w{=}25.7$ mV, solid)\\$328$ K ($w{=}28.3$ mV, dotted-bold)};
% ============ bottom panel: dxi_eq/dV bell at same three temperatures ============
\begin{scope}[yshift=-5.4cm,y=0.255cm]
\draw[-{Latex[length=1.6mm]}] (-108,0) -- (130,0) node[below,font=\scriptsize] {$V-U_j^{\,d}$ [mV]};
\draw[-{Latex[length=1.6mm]}] (-108,0) -- (-108,11.7) node[above,font=\scriptsize,align=center] {$\dd\xi_\eq/\dd V$\\{\tiny$[\mathrm{V}^{-1}]$}};
\foreach \xx in {-100,-50,50,100} {\draw (\xx,0.14) -- (\xx,-0.14) node[below,font=\scriptsize] {\xx};}
\foreach \yy in {4,8} {\draw (-104,\yy)--(-112,\yy) node[left,font=\scriptsize] {\yy};}
\draw[densely dotted,black!45] (0,0) -- (0,11.2);
\draw[densely dashed] plot[smooth] coordinates {(-100,0.5554) (-90,0.8444) (-80,1.2744) (-70,1.9019) (-60,2.7915) (-50,3.9984) (-40,5.5308) (-30,7.2915) (-20,9.0246) (-10,10.3331) (0,10.8251) (10,10.3331) (20,9.0246) (30,7.2915) (40,5.5308) (50,3.9984) (60,2.7915) (70,1.9019) (80,1.2744) (90,0.8444) (100,0.5554)};
\node[font=\tiny] at (0,11.45) {$10.83$ ($268$K)};
\draw[thick] plot[smooth] coordinates {(-100,0.7616) (-90,1.1031) (-80,1.584) (-70,2.2463) (-60,3.13) (-50,4.2555) (-40,5.5964) (-30,7.0453) (-20,8.3964) (-10,9.3754) (0,9.7353) (10,9.3754) (20,8.3964) (30,7.0453) (40,5.5964) (50,4.2555) (60,3.13) (70,2.2463) (80,1.584) (90,1.1031) (100,0.7616)};
\node[font=\tiny] at (32,9.0) {$9.74$ ($298$K)};
\draw[densely dotted,very thick] plot[smooth] coordinates {(-100,0.9713) (-90,1.351) (-80,1.861) (-70,2.5299) (-60,3.3778) (-50,4.403) (-40,5.5627) (-30,6.7565) (-20,7.824) (-10,8.5738) (0,8.8449) (10,8.5738) (20,7.824) (30,6.7565) (40,5.5627) (50,4.403) (60,3.3778) (70,2.5299) (80,1.861) (90,1.351) (100,0.9713)};
\node[font=\tiny] at (55,7.0) {$8.84$ ($328$K)};
\node[font=\scriptsize] at (0,-2.0) {center $U_j^{\,d}$: peak height $=$ central slope $=1/(4w)$};
\end{scope}
\end{tikzpicture}
\caption{평형 진행률과 그 미분, 온도 의존성(식~\eqref{eq:xieq}, $w_j=n_jRT/F$, $n_j{=}1$ — 좌표는 식 그대로의
수치 평가, $268/298/328$ K 세 온도). 위: $\xi_\eq(V)$ logistic — 낮은 $T$(좁은 $w$)일수록 더 가파르게 전환된다.
아래: 미분 $\dd\xi_\eq/\dd V=\xi_\eq(1-\xi_\eq)/w$ — 중심 $V=U_j^{\,d}$ 에서 최대이며, 그 값이 곧 중심 기울기
$1/(4w)$(폭의 기하적 의미) — $268$ K 에서 $10.83\,\mathrm V^{-1}$, $298$ K 에서 $9.74\,\mathrm V^{-1}$, $328$ K
에서 $8.84\,\mathrm V^{-1}$ 로 온도가 오를수록 폭이 넓어져 peak 가 낮아진다. 이 종이 다음 절의 평형 peak 이고,
방향에 불변(양수)이다.}
\label{fig:logistic}
